\documentclass{article}

\usepackage{enumerate}
\usepackage{hyperref}
\usepackage{amsmath}
% \usepackage{amsfonts}

\usepackage{xcolor}
\newcommand{\jarad}[1]{{\color{blue} Jarad: #1}}
\newcommand{\spencerr}[1]{{\color{red} Response: #1}}
\newcommand{\spencerq}[1]{{\color{green} Question: #1}}

\begin{document}
  
  
\Huge
AE1
\normalsize
\section{Main Comments}
\begin{enumerate}[(1)]
  \item Having read the paper by Loaiza-Maya et al. (2021) when it was 
  initially published, it
appears to me that the approach suggested in (4) is a specific instantiation 
of the “Focused
Bayesian predictive” proposed by those authors in the specific case of linear 
pools. I also note
that such a case was indeed considered by those authors in that paper: in 
their Section 4.1,
Loaiza-Maya et al. (2021) directly apply the Gibbs posterior framework for a 
mixture of two
time series models used for financial returns (see their Equation (17) on page 
13 in the arXiv
version). The direct link between the current proposal and this earlier work 
is not clear enough
in the current draft and must be clarified right after the stacking posterior 
is introduced.

\spencerq{Things that I believe right off the bat that make this new
\begin{itemize}
  \item I provide a new definition of the CRPS which can accommodate a mixture
    of continuous distributions for any finite number of components
  \item The SGP works for any number of components, not just 2 and not only for 
      normal components
  \item The decay rate is added to my loss function
\end{itemize}
}


  \item On page 6, it is odd on the one hand to write in (5) that
  \[
    \theta^* \in \mathrm{arg min} R(\theta)
  \]
  but then for the sample optimisation problem to write
  \[
    \hat{\theta}_n  = \mathrm{arg min} R_n(\theta)
  \]
We reserve the notation $\in$ to clarify that the minimiser is not unique. 
But if the population
minimiser is not unique, then why would the sample minimiser be unique?

\item  On page 7, when you talk about the selection of the learning rate, it 
may be useful to
point out that there is a meaningful sense in which predictions become ``more`` 
or ``less`` accurate for different choices of the learning rate. This fact was 
indeed suggested in Loaiza-Maya et
al. (2021) but was recently rigorously proven to be a valid conjecture by 
McLatchie et al. (2025).

\item I found the use of the two different loss functions based on the 
dependence in the data to
be a useful addition to the literature. The construction itself appears very 
similar to the “Loss
Discounting Framework” proposed in Section 3.1 of Bernaciak and Griffin (2024).
Of course,
you are discounting within the loss used to construct the posterior and they 
are discounting the
predictive itself. Still, this should be mentioned.

\item  In Section 3.3, why do we care about consistency here? There is little 
to no motivation here
and that should be rectified. One way to argue for this is that posterior 
consistency can be
used to show that the resulting posterior predictive is more accurate than that 
obtained via the
Bayes posterior. When you talk about posterior consistency, you should bolster 
your argument by mentioning that Loaiza-Maya et al. (2021) and Frazier et al. 
(2024) have already obtained
such results under certain conditions.

\item by mentioning that Loaiza-Maya et al. (2021) and Frazier et al. 
(2024) have already obtained such results under certain conditions.

  \begin{enumerate}[1.]
    \item On page 28 you talk about the DP prior, but the prior condition is 
    stated as general in
  your result. Rather than stating this, I suggest showing that this condition 
  is satisfied in
  your example by the DP prior and giving it as a separate result in the 
  appendix.
  
  \spencerq{This seems weird to me. Why would a more specific case be desired
            if it's covered under the given condition?}
    \item Why have you only considered the CRPS in this result? My guess is that 
    this is because
    you can not generally find a simple upper bounding function for the log 
    score; 
    however,
    Jensen should provide such a condition since $-\text{log}$ is concave.
  \end{enumerate}
  
\item Page 11, for the learning rate in the example of Section 4.1, what sample 
size was used for
the selection of $\eta$ when using the CRPS? Please clarify in the supplement 
precisely how this
optimisation was performed.

\item In Section 4.3, how was the posterior sampling conducted? In particular, 
was the CRPS for
this mixture actually available in closed-form or did an estimator need to be 
used within the
posterior? Given that one of the models was an SIR model, I would be surprised 
if the predictive
was Gaussian, which would mean getting the CRPS in closed-form would not be 
feasible.


\spencerr{Clarified by adding:
``Neither of the two Bayesian forecast component models produced 
closed form
CDFs, but rather the distributions
were estimated from 10,000 draws from the posterior
predictive distributions. Draws were also sampled from the the ARIMA and random
walk predictive distributions, and to fit the SGP 
the CRPS was approximated using equation
(15).``}

\item The SGP50 and the EQW have quite similar histograms. What does the 
posterior for the
SGP50 look like? Is it just basically equal weight across the models?


\end{enumerate}

\subsection*{Minor Comments}
\begin{enumerate}[1.]
  \item Page 4, ``$h$ time steps into the with``. Missing ``future`` I guess.
  
  \spencerr{added ``future`` to make sentence complete}
  
  \item Page 6, when you introduce $R(\theta)$ as $El_{\theta}(y)$ 
  clarify that $E$ is 
  expectation wrt the true data
generating distribution of the data $y$.

\spencerr{added, ``where $E$ denotes the expected value with 
respect to the true data generating distribution of $y_i$.``}

  \item Page 19, ``a thus a stronger``
  
  \spencerr{fixed and edited to now read ``thus provides stronger``}
  \item The grey scale in Figure 4 cannot be easily interpreted. Please 
  use a different colour scheme
or way of denoting the methods.
\end{enumerate}
\newpage






\Huge
R1
\normalsize
\\~\\
In this paper, the authors develop a generalized Bayesian estimation framework 
for determining the 
weights of a mixture of given predictive distributions. Unlike standard 
Bayesian updating based on the 
likelihood function, their approach employs proper scoring rules, leading to 
posterior distributions 
over the simplex that concentrate on weights delivering high predictive 
accuracy under the chosen 
score. The paper illustrates the framework through a series of numerical 
experiments and an 
application to the influenza forecasting competition.
\\~\\
The paper is clearly written and thorough, but I view it more as an in-depth 
elaboration of ideas 
already present in the literature. For example, Section 4.1 of Loaiza-Maya et 
al. (2021) (hereafter LMF) 
also considers mixtures of given predictive distributions, with mixture weights 
estimated via 
generalized Bayesian inference under various scoring rules, including the CRPS. 
The main difference is 
that LMF restricted attention to two-component mixtures. Moreover, Lemma 1 in 
LMF establishes a 
concentration result (which extends naturally to mixtures with more components), 
closely paralleling 
the implications of Theorem 1 in the present paper. 
\\~\\
Another relevant connection is Frazier et al. (2024) (hereafter FLMK). In 
Example 5.2, FLMK study 
mixtures of normals with time-varying weights, where a Gibbs posterior—updated 
via proper scoring 
rules—is placed on all model parameters, including both weights and mixture 
components. While 
FLMK advocate variational inference for computational tractability, the target 
remains the exact Gibbs 
posterior defined by the scoring rule.
\\~\\

\spencerq{The work by Frazier is relevant in that it produces a Gibbs posterior
on model parameters, but the difference in what I am doing is that the models
are already given to me. I am not estimating model parameters. I am only 
estimating the weights of a convex combination of distribution functions from
the various given models. They also say that the CRPS cannot be calculated in 
the case of the mixture }
The main novel aspect I see in this paper is the use of a discounting factor 
in the loss function to deal 
with the potential time varying nature of the data generating process. 
However, I do not believe this 
aspect to have been developed or emphasised enough.
\\~\\
Taken together, these connections suggest that although this paper is 
carefully executed and well 
written, its contribution lies primarily in consolidating and extending 
previously introduced ideas.


\end{document}